\chapter{历法}
\chapterkana{暦}
\epigraphbox{“乃命羲和，钦若昊天，历象日月星辰，敬授人时。”}{——《尚书·尧典》}

\kaishou{历法管三件事：下种、动土、杀人。\par
谁定历，谁就定这三件事的日子。\par
后羿改过一次历。此后夏的天，在他手里。}

\section{规程里的判据}

前面几章用的材料，都是维持性的：住哪里、怎么进出、吃什么、谁在册。这套东西管的是让城活下去，两万年不需要改。

城里另有一套文字，管的是外面，与维持无关。基金会编目时把它们合称《规程残卷》，编号 \texttt{ФВ-1971-22}，共存十一叶，叶次不连续，最大的一叶编号到四十四。

第四十四叶只有三行，中间一行是一句设问，抄手写的那一笔比上下文都用力：

\begin{archive}[内部参考]{档案 № ФВ-1971-22}
第四十四叶，共三行。上行：「凡遣物于地，先定其判据。」中行：「土著得历之后，作何状。」下行：「记之，勿断。」
叶末另有一行小字，墨色较淡：「此条自登船前立，未改。」
\end{archive}

三行字里，要紧的是「未改」两个字。前面说过，章程总目一千二百四十七条，修订栏全空；这条规程不在总目里，却带着同样一句话：自登船前立，未改。

为什么没改，第二章给过解释：改一条要重新判断轻重，而轻重已经冻住。这个解释在这里比在总目那里更难反驳，因为这一条不是维持性的条文，它是一句设问。设问不解决任何问题，划掉它也不会让城停转。两万年里没有人划掉它，只能说明没有人觉得它值得一划。

\section{投放清单}

规程残卷里另有一叶，把要送下去的东西分成两栏，一栏可遣，一栏不可遣。

可遣一栏短：一部历法，青铜配方，灌溉法。不可遣一栏长：文字，高阶冶炼，武器，历法的完整算法。

清单只到目录，没有正文。哪一部历、配方几成、灌溉法是哪一种，残卷里都没有。能补的只有一处旁证：中原的观象遗址。

山西襄汾的陶寺遗址，一九七八年起发掘，城址的年代约在公元前两千一百年至前一千九百年之间。遗址东南有一组半圆形夯土柱，共十三根，柱间形成十二道缝；站在既定的一点上，从缝里看日出，一年之中只有少数几天的日出会正对某一道缝。这是一台用建筑做成的历。\footnote{陶寺观象台复原方案由考古队与中国科学院自然科学史研究所合作提出，十三柱十二缝的观测结论发表于二〇〇〇年前后；此后有不同意见，认为部分柱基未必属于同一期。}

它的年代，比劣化版历法送到中原要早一百年。这一点很要紧：中原不是没有能力观天，它有。所以「不授高阶算法」这一条，删掉的不是能力，而是精确度。中原本来就会观天，只是观得粗。

至于文字为什么不给，规程没有写理由。笔者能想到的解释只有一个：文字会累积。有文字的地方，上一代算错的东西，下一代能看见；没有文字的地方，每一代都从头算。这句话没有材料作证，只是笔者自己想的。

\section{十部与一部}

劣化版历法送到中原，约在公元前二千年。这一节往后所讲的，大部分要靠地面文献，而地面文献说的是神话。

先说清材料的分量。《左传》襄公四年记后羿「因夏民以代夏政」，这是史书；《淮南子·本经训》记「十日并出」，这是神话。两件事在书里是两件事，在档案里没有对应物。

把两边合起来看，可以拼出一条线。历法送到之后，夏把它抄成若干部，分授各方国。有穷氏是其中之一，其首领后羿手上那一部，他用了半辈子。

公元前二千年，后羿代夏。他做的第一件事是废历：其余各部一律收回销毁，只留自己手上那一部。

《淮南子》把这件事记成「十日并出，羿射九日」。十日并出，说的是很多部历同时行用；羿射九日，说的是废掉九部。这个读法有一个好处：它解释了神话里一个多余的数字。天上多出九个太阳，对农业毫无意义；十个授时中心同时存在，对权力很有意义。

后羿没有把十部合并成一部，他选了最费事的那条路：全部收回，全部销毁。合并要把十部收齐、逐部校对、重定历元，他不识字，做不了；收回与销毁只需要一道命令。

这一节最弱的地方在这里：十这个数字，是从神话里来的。档案没有说夏抄了几部，规程残卷也没有。笔者用它，是因为它恰好能解释神话里那个多余的九，而别的解释做不到这一点。这是猜，猜得通不证明猜得对。

\section{信道与底牌}

劣化版历法删掉的三样，岁差、交食、行星逆行，都是算得出来却看不见的东西。删掉它们，足够定农时，不足以看天。

但清单里还有一条，规程写得很轻：「余一段，无用，留之。」留下来的这一段，就是后羿后来补回去的那一段。

后羿补历用了五年。他不识字，靠的是观测：冬至对不上，交食没有如期发生。五年之后，他把缺的算法逐项填回，自称「新历」。他以为那是一部修好的历法。

修好之后，他手上多了一样东西：一套能处理天体的算法。他一生没有用过，因为他不知道那些数是干什么的。

历法还有第二层，这一层是月之民自己留的，后羿从未发觉。历法条文里夹着一条信道，方向双向。上行是人的作息：按历下种，按历动土，按历杀人，作息本身即是回传。下行是条文：月之民可以把话写成历法用语，混在农时、祭祀、刑杀的规定里，下面的人按历行事，等于照办。

后羿不识字，分不出哪一条是外来的。用今天的话说，他手上那部历既是日历，也是收报机和发报机，而他不认得报文。

公元前一九九〇年前后，新历的条文经这条信道上行，送回月面。月之民一眼认了出来：自己的底牌，出现在一个土著手里。

从这一刻起，地面第一次出现了能威胁月球的东西。持有者既不知道，也不会用，但它在那儿。此事之后不到十年，月面开始收场，那一段留到第五章。

这一章用了三样东西：规程残卷、地面神话、陶寺的柱子。柱子和年代能复核，残卷只有十一叶，神话最不可靠。笔者把神话摆在论证位置上，是这一章最冒险的一处；若哪天有人从别的材料里读出夏只抄了五部，这个读法就要重排。

笔者不打算把话说得更圆。神话能解释的东西，往往不止一种解释；能用它，是因为别处没有更好的材料，与它可靠与否无关。
